\title{The Era of 1-bit LLMs: All Large Language Models are in 1.58 Bits}

\begin{abstract}
Recent advancements, exemplified by BitNet~\cite{bitnet}, signal the onset of a new chapter in the evolution of 1-bit Large Language Models (LLMs). This paper presents \textbf{\text{BitNet b1.58}}, a 1-bit LLM implementation in which each parameter (or weight) is ternary, taking values from the set \{-1, 0, 1\}. This model achieves parity with its full-precision (FP16 or BF16) Transformer LLM counterparts of identical model size and number of training tokens, matching both perplexity and end-task results. At the same time, it offers considerable advantages in terms of efficiency, manifesting in lower latency, reduced memory footprint, increased throughput, and diminished energy usage. Notably, the 1.58-bit LLM introduces a fresh perspective on scaling laws and offers a novel framework for constructing future LLMs that optimize both performance and cost. This innovation not only suggests the possibility for specialized computational paradigms but also paves the way for designing hardware tailored specifically to 1-bit LLMs.
\end{abstract}

\section{The Era of 1-bit LLMs}

Over recent years, Artificial Intelligence has witnessed unprecedented advancements in the development and capabilities of Large Language Models (LLMs). These models have excelled across diverse natural language processing applications, yet their escalating parameter counts have led to significant practical challenges, particularly concerning deployment and growing worries over their environmental footprint and associated financial costs stemming from increased energy consumption.

A prominent solution to mitigate these issues has involved post-training quantization, whereby models are converted to operate with low-precision weights and activations for inference~\cite{smoothquant, gptq, quip, quip_sharp}. Such approaches yield dramatic savings in both memory usage and computational complexity. Progressively, the field has transitioned from 16-bit representations to lower bit-widths, including 4-bit versions~\cite{gptq, awq}. Despite widespread adoption within the industry, post-training quantization methods have inherent limitations and do not represent the optimal strategy.

Recent advancements in 1-bit neural network architectures, exemplified by BitNet~\cite{bitnet}, suggest an effective route for minimizing the resource requirements of LLMs without sacrificing accuracy. Standard LLMs employ 16-bit floating-point numbers (FP16 or BF16), with matrix multiplications constituting the core computational workload. The primary computational burden stems from the floating-point addition and multiplication these operations entail. By contrast, BitNet restricts matrix multiplications to integer additions, drastically reducing the energy demands for LLM operation. Given that many hardware platforms are limited by power consumption, such energy reductions can translate into accelerated computational speed.

Beyond raw computation, the transfer of model parameters from DRAM to on-chip accelerator memory (such as SRAM) also constitutes a major cost during inference. Strategies have been proposed to expand SRAM in order to boost throughput; however, this comes with substantially increased costs relative to DRAM. 1-bit LLMs possess a far smaller memory footprint than their full-precision counterparts, in both storage requirements and bandwidth consumption. Consequently, this leads to lower cost and reduced latency in loading weights from DRAM, facilitating swifter and more efficient inference processes.

In this paper, we present a notable 1-bit LLM extension, \textbf{\text{BitNet b1.58}}, in which all parameters are ternary, assuming values from the set \{-1, 0, 1\}. By introducing 0 as an additional discrete state to the original 1-bit BitNet, we obtain a representation equivalent to 1.58 bits in binary terms. \text{BitNet b1.58} inherits the computational advantages of its predecessor, such as a fundamentally different operation paradigm for matrix multiplication that nearly eliminates multiplication and supports substantial computational optimization. The energy profile remains on par with the original BitNet, and the model exhibits significant improvements in memory efficiency, throughput, and inference latency relative to FP16-based LLMs. In addition, \text{BitNet b1.58} brings two further improvements. The first is an enhanced representational power, made possible by the explicit integration of 0 as a model parameter, which enables feature filtering and notably boosts the performance potential of 1-bit LLMs. The second is empirical: our results demonstrate that, beginning at a 3B parameter scale and under matching conditions (including model size and training token count), \text{BitNet b1.58} is able to equal the perplexity and downstream task performance of full-precision (FP16) models.

\section{BitNet b1.58}

\text{BitNet b1.58} extends the BitNet model, a variant of the Transformer architecture that substitutes the standard \emph{nn.Linear} layers with \emph{BitLinear} layers. This model is trained entirely from scratch, employing 1.58-bit quantization for weights and 8-bit quantization for activations. In comparison to the initial BitNet, several alterations have been made, which are detailed below.

\paragraph{Quantization Function.} To ensure that the weights assume values exclusively from the set $\{-1, 0, +1\}$, we use the \emph{absmean} quantization approach. This method normalizes the weight matrix using its mean absolute value, followed by rounding each normalized element to the nearest integer in $\{-1, 0, +1\}$ as follows:

\begin{equation}
    \widetilde{W} = \mathrm{RoundClip}\left(\frac{W}{\gamma + \epsilon}, -1, 1\right),
\end{equation}

\begin{equation}
    \mathrm{RoundClip}(x, a, b) = \max(a, \min(b, \mathrm{round}(x))),
\end{equation}

\begin{equation}
    \gamma = \frac{1}{nm}\sum_{ij} |W_{ij}|.
\end{equation}

For activations, the quantization scheme remains consistent with that of BitNet, except that we omit the rescaling of activations to the interval $[0, Q_b]$ before applying nonlinearities. Instead, activations are rescaled on a per-token basis to $[-Q_b, Q_b]$, thus removing zero-point quantization. This modification offers practical advantages in implementation and system optimization, while our experiments demonstrate negligible impact on overall model accuracy.

\paragraph{LLaMA-alike Components.}

LLaMA~\cite{llama, llama2} serves as a foundational architecture for a large fraction of open-source LLMs. To maximize compatibility with open-source frameworks, \text{BitNet b1.58} incorporates LLaMA-like elements. The model specifically adopts RMSNorm~\cite{rmsnorm}, SwiGLU~\cite{swiglu}, rotary embeddings~\cite{rope}, and omits all bias terms. As a result, \text{BitNet b1.58} is straightforward to use within popular open-source toolkits, such as Huggingface, vLLM~\cite{vllm}, and llama.cpp\footnote{\url{https://github.com/ggerganov/llama.cpp}}, facilitating seamless integration.

\section{Results}

We conducted a comprehensive comparison between \text{BitNet b1.58} and our reproduced FP16 \text{LLaMA LLM} models at various scales. Both models were pre-trained for 100 billion tokens using the RedPajama dataset~\cite{redpajama}, providing a consistent baseline for assessment. We assessed zero-shot capabilities across an array of language benchmarks, such as ARC-Easy~\cite{arc}, ARC-Challenge~\cite{arc}, Hellaswag~\cite{hellaswag}, Winogrande~\cite{winoGrande}, PIQA~\cite{piqa}, OpenbookQA~\cite{openbookqa}, and BoolQ~\cite{boolq}. Additionally, validation perplexities were computed on both the WikiText2~\cite{wikitext2} and C4~\cite{c4} corpora.

To analyze resource efficiency, we measured the runtime GPU memory footprint and inference latency for \text{LLaMA LLM} and \text{BitNet b1.58}. These metrics were recorded using the FasterTransformer\footnote{\url{https://github.com/NVIDIA/FasterTransformer}} library, recognized for its GPU inference optimization. For \text{BitNet b1.58}, the integration of the 2-bit kernel from Ladder~\cite{ladder} was employed. We focused on reporting the average time per output token, as it constitutes the principal cost during inference.

Table~\ref{tab:ppl} provides an overview of both perplexity and computational resource usage for \text{BitNet b1.58} and \text{LLaMA LLM}. The data illustrate that at the 3B scale, \text{BitNet b1.58} achieves perplexity comparable to the full-precision \text{LLaMA LLM}, while operating at 2.71$\times$ increased speed and consuming 3.55$\times$ less GPU memory. Notably, the 3.9B configuration of \text{BitNet b1.58} not only is 2.4$\times$ faster and 3.32$\times$ more memory efficient, but also exhibits substantial performance gains relative to the \text{LLaMA LLM} 3B counterpart.

In Table~\ref{tab:zero-shot}, we detail the zero-shot accuracy metrics for various end tasks. Evaluation procedures followed the \emph{lm-evaluation-harness}\footnote{\url{https://github.com/EleutherAI/lm-evaluation-harness}} framework. The findings indicate a diminishing performance gap between \text{BitNet b1.58} and \text{LLaMA LLM} with increasing model size. Importantly, from the 3B scale onward, \text{BitNet b1.58} attains parity with its full precision equivalent. Echoing the perplexity trends, the 3.9B variant of \text{BitNet b1.58} consistently surpasses the 3B \text{LLaMA LLM} model in both effectiveness and efficiency. Collectively, these results demonstrate that \text{BitNet b1.58} constitutes a Pareto improvement over existing state-of-the-art large language models.

\paragraph{Memory and Latency}

The model was further expanded to 7B, 13B, and 70B parameters, and corresponding cost assessments were carried out. As depicted in Figure~\ref{fig:latency-memory-energy}, both latency and memory usage exhibit clear scaling patterns, with greater speed improvements observed as the model size increases. Notably, \text{BitNet b1.58} at 70B achieves a speed-up of 4.1 times over the \text{LLaMA LLM} baseline. This improvement arises from the increased time demands of the \emph{nn.Linear} layer as models become larger. Memory usage also escalates accordingly, since the embeddings retain full precision, resulting in a smaller proportional memory footprint for larger models. All measurements for latency and memory relied on a 2-bit kernel implementation, indicating the potential for further cost reduction through additional kernel optimizations.

\paragraph{Energy}

We additionally approximated the energy required for arithmetic operations in both \text{BitNet b1.58} and \text{LLaMA LLM}. The primary focus was on the energy associated with matrix multiplication, as it dominates the resource expenditure in large language models. Figure~\ref{fig:energy} presents the breakdown of these energy costs. For \text{BitNet b1.58}, most computations are INT8 additions, whereas for \text{LLaMA LLM} both FP16 additions and FP16 multiplications are prominent. Based on the energy estimation framework provided in~\cite{energycost, pokebnn}, matrix multiplication with \text{BitNet b1.58} yields a 71.4-fold reduction in arithmetic operation energy consumption compared to standard methods on 7nm hardware. Further, we analyzed total energy requirements for end-to-end processing on sequences of 512 tokens. Our findings demonstrate that \text{BitNet b1.58} offers progressively greater energy efficiency relative to the FP16 \text{LLaMA LLM} as model size grows. This trend can be attributed to the increasing dominance of the \emph{nn.Linear} operation in larger models and the relatively diminishing contributions from other architectural elements.

\paragraph{Throughput}

We assessed the throughput of both \text{BitNet b1.58} and \text{LLaMA LLM} with 70B parameters utilizing two 80GB A100 GPUs, employing pipeline parallelism~\cite{gpipe} to enable execution of \text{LLaMA LLM} 70B on this hardware. By incrementally increasing the batch size until GPU memory was exhausted, with the sequence length set to 512, we observed the performance outlined in Table~\ref{tab:throughput}. Notably, \text{BitNet b1.58} 70B accommodates a batch size as much as 11 times greater than \text{LLaMA LLM}, yielding a throughput increase by a factor of 8.9.

\textbf{\text{BitNet b1.58} introduces a new paradigm in the scaling relationship between model quality and inference resource requirements}. Drawing from the data in Figure~\ref{fig:latency-memory-energy} and \ref{fig:energy}, we note the following model size equivalences between the 1.58-bit and conventional 16-bit settings:

\begin{itemize}
    \item 13B BitNet b1.58 surpasses the 3B FP16 LLM with respect to latency, memory footprint, and power usage.
    \item 30B BitNet b1.58 offers better efficiency in latency, memory consumption, and energy use compared to 7B FP16 LLM.
    \item 70B BitNet b1.58 outperforms the 13B FP16 LLM across latency, memory requirements, and energy demands.
\end{itemize}

\paragraph{Training with 2T Tokens}

The amount of training data, quantified in tokens, plays a pivotal role in the effectiveness of LLMs. To evaluate the token scalability of \text{BitNet b1.58}, we conducted experiments by training a \text{BitNet b1.58} model on 2T tokens, adopting the same data methodology as StableLM-3B~\cite{StableLM-3B-4E1T}—currently a leading open-source 3B model. Performance was assessed using a benchmark suite consisting of Winogrande~\cite{winoGrande}, PIQA~\cite{piqa}, SciQ~\cite{sciq}, LAMBADA~\cite{lambada}, and ARC-easy~\cite{arc}, with zero-shot accuracies provided in Table~\ref{tab:2t}. When tasks involved both standard and normalized accuracy measures, their average was reported. The outcomes for StableLM 3b at the 2T token level were cited from its technical documentation. The evaluation shows that \text{BitNet b1.58} consistently achieves higher scores on every evaluated benchmark, suggesting that 1.58-bit LLMs maintain robust generalization performance.

\section{Discussion and Future Work}

\textbf{1-bit Mixture-of-Experts (MoE) LLMs}

The Mixture-of-Experts (MoE) strategy has established itself as an efficient solution for large language models by significantly lowering computational FLOPs. However, substantial memory requirements and extensive inter-chip data transfer remain barriers to practical use. Adopting 1.58-bit LLMs has the potential to overcome these issues. By minimizing memory usage, fewer hardware units are necessary for MoE deployment, and communication costs for moving activations across devices are markedly reduced. If the entire model can reside on a single chip, network transfer bottlenecks could be entirely eliminated.

\textbf{Native Support of Long Sequence in LLMs}

Supporting long input sequences natively in large language models has emerged as a vital need. One of the primary obstacles is the significant memory required to maintain KV caches during inference with extended contexts. \text{BitNet b1.58} offers a meaningful advancement in this area by lowering activation precision from 16 bits to 8 bits, which, in turn, permits twice the context length within the same hardware limits. Prospects for further work include the lossless compression of activations to 4 bits or even smaller representation in 1.58-bit LLMs.

\textbf{LLMs on Edge and Mobile}

Deploying 1.58-bit LLMs can substantially enhance language modeling capabilities on edge and mobile platforms, which typically face strict constraints on memory and compute resources. With their lower demands on both memory and power, 1.58-bit LLMs make possible a range of previously unfeasible applications in these environments. Such advances may significantly expand the utility and application scope of LLMs on resource-constrained devices. Additionally, as CPUs remain the predominant processing units on edge and mobile hardware, the compatibility of \text{BitNet b1.58} with these processors allows for even more efficient operation and wider practical deployment.

\textbf{New Hardware for 1-bit LLMs}

Innovations such as Groq\footnote{\url{https://groq.com/}} have showcased the effectiveness of purpose-built hardware, such as LPUs, tailored for large language models. Building upon this trend, there is an opportunity—and a need—to develop hardware and system designs specifically optimized for the unique requirements of 1-bit LLMs, as introduced in BitNet~\cite{bitnet}, to fully leverage this new computational approach.

\end{document}